\documentclass[../../../include/open-logic-section]{subfiles}

\begin{document}
\olfileid{sth}{card-arithmetic}{expotough}

\olsection[কিছু সরলীকরণ]{অঙ্কবাচক সংখ্যার ঘাতে কিছু সরলীকরণ}

$\canonord$ সংজ্ঞায়িত করা একটু জটিল ছিল। কিন্তু তার
ফলে অঙ্কবাচক সংখ্যার যোগ ও গুণের উপযোগী ফল
\olref[simp]{cardplustimesmax} পেয়েছি। ট্রান্সফাইনাইট
ঘাতকে অবশ্য এত সহজে সরল করা যায় না। কেন যায় না,
তা বোঝার আগে সসীম ক্ষেত্রের পরিচিত ধরনটি প্রসারিত
করে একটি ফল দেখি (যদিও তার প্রমাণ বেশ বিমূর্ত):

\begin{prop}\ollabel{simplecardexpo}
যেকোনো অঙ্কবাচক সংখ্যা $\cardfont{a},\cardfont{b},\cardfont{c}$-র জন্য
$\cardexpo{\cardfont{a}}{\cardfont{b}\cardplus\cardfont{c}}=
\cardexpo{\cardfont{a}}{\cardfont{b}}\cardtimes
\cardexpo{\cardfont{a}}{\cardfont{c}}$ এবং
$\cardexpo{(\cardexpo{\cardfont{a}}{\cardfont{b}})}{\cardfont{c}}=
\cardexpo{\cardfont{a}}{\cardfont{b}\cardtimes\cardfont{c}}$।
\end{prop}

\begin{proof}
প্রথম দাবির জন্য
$f\colon(\cardfont{b}\disjointsum\cardfont{c})\to\cardfont{a}$
ধরি। একে ``ভাগ'' করতে প্রত্যেক
$\beta\in\cardfont{b}$-র জন্য
$f_\cardfont{b}(\beta)=f(\beta,0)$, এবং প্রত্যেক
$\gamma\in\cardfont{c}$-র জন্য
$f_\cardfont{c}(\gamma)=f(\gamma,1)$
সংজ্ঞায়িত করি।
% BN-SRC-456: an element of the product of function spaces is a pair of functions.
$f\mapsto\tuple{f_\cardfont{b},f_\cardfont{c}}$
চিত্রণটি একটি !!a{bijection}
$\funfromto{\cardfont{b}\disjointsum\cardfont{c}}{\cardfont{a}}
\to(\funfromto{\cardfont{b}}{\cardfont{a}}\times
\funfromto{\cardfont{c}}{\cardfont{a}})$।

দ্বিতীয় দাবির জন্য
$f\colon\cardfont{c}\to
(\funfromto{\cardfont{b}}{\cardfont{a}})$
ধরি। তাহলে প্রত্যেক $\gamma\in\cardfont{c}$-র জন্য
$f(\gamma)\colon\cardfont{b}\to\cardfont{a}$
একটি অপেক্ষক। প্রত্যেক
$\tuple{\beta,\gamma}\in\cardfont{b}\times\cardfont{c}$-র
জন্য $f^*(\beta,\gamma)=(f(\gamma))(\beta)$
সংজ্ঞায়িত করি।
% BN-SRC-457: f* is literally a function on the Cartesian product;
% transport across its cardinal-equivalence supplies the claimed exponent law.
$f\mapsto f^*$ চিত্রণটি একটি !!a{bijection}
$\funfromto{\cardfont{c}}{(\funfromto{\cardfont{b}}{\cardfont{a}})}
\to\funfromto{\cardfont{b}\times\cardfont{c}}{\cardfont{a}}$।
আর কার্তেসীয় গুণফলটি তার অঙ্কবাচক গুণফলের
সমসংখ্যক;
একটি দ্বিপারস্পরিক চিত্রণ দিয়ে অপেক্ষকের
সংজ্ঞাক্ষেত্র পরিবহন করলেই দাবিকৃত ঘাত-সমতা মেলে।
\end{proof}

অসীম অঙ্কবাচক সংখ্যার জন্য
$\cardexpo{\cardfont{a}}{\cardfont{b}}$
সহজে হিসাব করার উপায় আমরা \emph{চাই}।
সেদিকে একটি ভালো পদক্ষেপ হলো:

\begin{prop}\ollabel{cardexpo2reduct}
$2\leq\cardfont{a}\leq\cardfont{b}$ এবং
$\cardfont{b}$ অসীম হলে
$\cardexpo{\cardfont{a}}{\cardfont{b}}=
\cardexpo{2}{\cardfont{b}}$।
\end{prop}

\begin{proof}
\begin{align*}
	\cardexpo{2}{\cardfont{b}} &\leq
	\cardexpo{\cardfont{a}}{\cardfont{b}}\text{, as $2 \leq \cardfont{a}$}\\
	&\leq \cardexpo{(2^\cardfont{a})}{\cardfont{b}}
	\text{, by \olref[opps]{lem:SizePowerset2Exp}}\\
	&= \cardexpo{2}{\cardfont{a} \cardtimes \cardfont{b}}
	\text{, by \olref{simplecardexpo}} \\
	&= \cardexpo{2}{\cardfont{b}}
	\text{, by \olref[simp]{cardplustimesmax}}
\end{align*}
\end{proof}

আর বেশি সরলীকরণ আশা করা ঠিক হবে না, কারণ
\olref[card-arithmetic][opps]{lem:SizePowerset2Exp}
অনুযায়ী $\cardfont{b}<\cardexpo{2}{\cardfont{b}}$।
তবে $\cardfont{b}<\cardfont{a}$ হলে
$\cardexpo{\cardfont{a}}{\cardfont{b}}$ নিয়ে
কী বলব, তা এই ফলে জানা যায় না। অবশ্য
$\cardfont{b}$ \emph{সসীম} হলে পথ জানা আছে।

\begin{prop}
% BN-SRC-458: zero exponent gives one, so the source proposition needs n nonzero.
$\cardfont{a}$ অসীম এবং $n\in\omega$ অশূন্য হলে
$\cardexpo{\cardfont{a}}{n}=\cardfont{a}$।
\end{prop}

\begin{proof}
$\cardexpo{\cardfont{a}}{n}=\cardfont{a}\cardtimes\cardfont{a}
\cardtimes\ldots\cardtimes\cardfont{a}=\cardfont{a}$,
\olref[simp]{cardplustimesmax} অনুযায়ী। শূন্য ঘাতে
অপেক্ষকের সংজ্ঞাক্ষেত্র শূন্য সেট, তাই ফল এক—উপরের
সমতা সেক্ষেত্রে প্রযোজ্য নয়।
\end{proof}
\noindent
আরও কিছু ক্ষেত্রে
$\cardexpo{\cardfont{a}}{\cardfont{b}}$-র
আকার নিয়ন্ত্রণ করা যায়:

\begin{prop}
$2\leq\cardfont{b}<\cardfont{a}\leq
\cardexpo{2}{\cardfont{b}}$ এবং $\cardfont{b}$
অসীম হলে
$\cardexpo{\cardfont{a}}{\cardfont{b}}=
\cardexpo{2}{\cardfont{b}}$।
\end{prop}

\begin{proof}
$\cardexpo{2}{\cardfont{b}}\leq\cardexpo{\cardfont{a}}{\cardfont{b}}
\leq\cardexpo{(\cardexpo{2}{\cardfont{b}})}{\cardfont{b}}=
\cardexpo{2}{\cardfont{b}\cardtimes\cardfont{b}}=
\cardexpo{2}{\cardfont{b}}$,
\olref{cardexpo2reduct}-এর মতো যুক্তিতে।
\end{proof}
\noindent
এর পরে বিষয়টি আরও সূক্ষ্ম হয়ে ওঠে।

\end{document}
